% GENERATED FILE. Edit canonical lessons, then run npm run generate:books.
% canonical-source-hash: fnv1a64:7039446941b3c3ee
% canonical-lessons: JA-C298-mukashi, JA-C298-mukashibanashi, JA-C298-kataru, JA-C298-mura, JA-C298-hatarakimono

\chapter{Long, Long Ago}
\label{ch:ja-long-long-ago}
\hlchaptermodality{298}

\begin{chapteropening}
\textbf{By the end of this chapter:} \emph{I can begin an old story and say who lived there.}

Complete the last lesson of chapter 298: I can begin an old story and say who lived there.
\end{chapteropening}

\section[\texorpdfstring{\ja{むかし} (mukashi)}{むかし (mukashi)}]{\ja{むかし} (mukashi) — long ago; the old days}

\label{lesson:JA-C298-mukashi}



\begin{quote}
Before the new one: say the Japanese for to drive out, then the Japanese for rare.
\end{quote}

\subsection*{You'll want to know: \ja{むかし}}
\textbf{\ja{むかし}} — \emph{mukashi} — \textquotedblleft{}long ago; the old days\textquotedblright{}.

\begin{quote}
\textbf{\ja{むかし、やまに} \ja{おじいさんが} \ja{すんでいました。}} — \emph{mukashi, yama ni ojiisan ga sunde imashita} — \textquotedblleft{}Long ago an old man lived in the mountains.\textquotedblright{}
\end{quote}

\begin{cousinweb}
\textbf{\ja{むかし}} (\emph{mukashi}) is a native Japanese word (wago), inherited rather than borrowed from Chinese; its kanji was attached later as a kun'yomi (JMdict).
\end{cousinweb}

\begin{grammarlens}[title={What you've built}]
One more word for telling what happened.
\end{grammarlens}

\subsection*{Your turn}
\begin{itemize}
\raggedright
  \item \emph{Say it:} \emph{mukashi}
  \item \emph{Say it:} \emph{mukashi}, once more
  \item \emph{Say it:} \emph{mukashi, yama ni ojiisan ga sunde imashita} — Long ago an old man lived in the mountains.
  \item \emph{Recall:} say the Japanese for to drive out, then the Japanese for rare, then say \emph{mukashi} again
\end{itemize}

\begin{tcolorbox}[breakable,colback=teal!4,colframe=teal!35!black,title={Before you move on}]
What does \ja{むかし} mean? (\textquotedblleft{}Long ago; the old days\textquotedblright{}.) Say it once more.
\end{tcolorbox}
\section[\texorpdfstring{\ja{むかしばなし} (mukashibanashi)}{むかしばなし (mukashibanashi)}]{\ja{むかしばなし} (mukashibanashi) — a folk tale; an old story}

\label{lesson:JA-C298-mukashibanashi}



\begin{quote}
Before the new one: say the Japanese for rare, then the Japanese for long ago.
\end{quote}

\subsection*{You'll want to know: \ja{むかしばなし}}
\textbf{\ja{むかしばなし}} — \emph{mukashibanashi} — \textquotedblleft{}a folk tale; an old story\textquotedblright{}.

\begin{quote}
\textbf{\ja{そぼは} \ja{いつも} \ja{むかしばなしを} \ja{しました。}} — \emph{sobo wa itsumo mukashibanashi o shimashita} — \textquotedblleft{}Grandmother was always telling old stories.\textquotedblright{}
\end{quote}

\begin{cousinweb}
\textbf{\ja{むかしばなし}} (\emph{mukashibanashi}) is a native Japanese word (wago), inherited rather than borrowed from Chinese, built from \emph{mukashi}, \textquotedblleft{}long ago\textquotedblright{}, and \emph{banashi}, \textquotedblleft{}talk\textquotedblright{} (from \emph{hanasu}) (JMdict, KANJIDIC2).
\end{cousinweb}

\begin{grammarlens}[title={What you've built}]
One more word for telling what happened.
\end{grammarlens}

\subsection*{Your turn}
\begin{itemize}
\raggedright
  \item \emph{Say it:} \emph{mukashibanashi}
  \item \emph{Say it:} \emph{mukashibanashi}, once more
  \item \emph{Say it:} \emph{sobo wa itsumo mukashibanashi o shimashita} — Grandmother was always telling old stories.
  \item \emph{Recall:} say the Japanese for rare, then the Japanese for long ago, then say \emph{mukashibanashi} again
\end{itemize}

\begin{tcolorbox}[breakable,colback=teal!4,colframe=teal!35!black,title={Before you move on}]
What does \ja{むかしばなし} mean? (\textquotedblleft{}A folk tale; an old story\textquotedblright{}.) Say it once more.
\end{tcolorbox}
\section[\texorpdfstring{\ja{かたる} (kataru)}{かたる (kataru)}]{\ja{かたる} (kataru) — to tell (a story); to relate}

\label{lesson:JA-C298-kataru}



\begin{quote}
Before the new one: say the Japanese for long ago, then the Japanese for a folk tale.
\end{quote}

\subsection*{You'll want to know: \ja{かたる}}
\textbf{\ja{かたる}} — \emph{kataru} — \textquotedblleft{}to tell (a story); to relate\textquotedblright{}.

\begin{quote}
\textbf{\ja{そふは} \ja{ときどき} \ja{むかしばなしを} \ja{かたります。}} — \emph{sofu wa tokidoki mukashibanashi o katarimasu} — \textquotedblleft{}Grandfather sometimes tells old stories.\textquotedblright{}
\end{quote}

Polite \textbf{\ja{かたります}}, \emph{katarimasu}; a godan verb. Past (polite): \textbf{\ja{そぼが} \ja{ゆめを} \ja{かたりました。}}, \emph{sobo ga yume o katarimashita}.

\begin{cousinweb}
\textbf{\ja{かたる}} (\emph{kataru}) is a native Japanese word (wago), inherited rather than borrowed from Chinese; its kanji was attached later as a kun'yomi (JMdict).
\end{cousinweb}

\begin{grammarlens}[title={What you've built}]
One more word for telling what happened.
\end{grammarlens}

\subsection*{Your turn}
\begin{itemize}
\raggedright
  \item \emph{Say it:} \emph{kataru}
  \item \emph{Say it:} \emph{kataru}, once more
  \item \emph{Say it:} \emph{sofu wa tokidoki mukashibanashi o katarimasu} — Grandfather sometimes tells old stories.
  \item \emph{Recall:} say the Japanese for long ago, then the Japanese for a folk tale, then say \emph{kataru} again
\end{itemize}

\begin{tcolorbox}[breakable,colback=teal!4,colframe=teal!35!black,title={Before you move on}]
What does \ja{かたる} mean? (\textquotedblleft{}To tell (a story); to relate\textquotedblright{}.) Say it once more.
\end{tcolorbox}
\section[\texorpdfstring{\ja{むら} (mura)}{むら (mura)}]{\ja{むら} (mura) — a village}

\label{lesson:JA-C298-mura}



\begin{quote}
Before the new one: say the Japanese for a folk tale, then the Japanese for to tell (a story).
\end{quote}

\subsection*{You'll want to know: \ja{むら}}
\textbf{\ja{むら}} — \emph{mura} — \textquotedblleft{}a village\textquotedblright{}.

\begin{quote}
\textbf{\ja{その} \ja{むらは} \ja{やまの} \ja{なかに} \ja{あります。}} — \emph{sono mura wa yama no naka ni arimasu} — \textquotedblleft{}The village lies in the mountains.\textquotedblright{}
\end{quote}

\begin{cousinweb}
\textbf{\ja{むら}} (\emph{mura}) is a native Japanese word (wago), inherited rather than borrowed from Chinese; its kanji was attached later as a kun'yomi (JMdict).
\end{cousinweb}

\begin{grammarlens}[title={What you've built}]
One more word for telling what happened.
\end{grammarlens}

\subsection*{Your turn}
\begin{itemize}
\raggedright
  \item \emph{Say it:} \emph{mura}
  \item \emph{Say it:} \emph{mura}, once more
  \item \emph{Say it:} \emph{sono mura wa yama no naka ni arimasu} — The village lies in the mountains.
  \item \emph{Recall:} say the Japanese for a folk tale, then the Japanese for to tell (a story), then say \emph{mura} again
\end{itemize}

\begin{tcolorbox}[breakable,colback=teal!4,colframe=teal!35!black,title={Before you move on}]
What does \ja{むら} mean? (\textquotedblleft{}A village\textquotedblright{}.) Say it once more.
\end{tcolorbox}
\section[\texorpdfstring{\ja{はたらきもの} (hatarakimono)}{はたらきもの (hatarakimono)}]{\ja{はたらきもの} (hatarakimono) — a hard worker}

\label{lesson:JA-C298-hatarakimono}



\begin{quote}
Before the new one: say the Japanese for to tell (a story), then the Japanese for a village.
\end{quote}

\subsection*{You'll want to know: \ja{はたらきもの}}
\textbf{\ja{はたらきもの}} — \emph{hatarakimono} — \textquotedblleft{}a hard worker\textquotedblright{}.

\begin{quote}
\textbf{\ja{むらの} \ja{わかい} \ja{おとこは} \ja{はたらきものでした。}} — \emph{mura no wakai otoko wa hatarakimono deshita} — \textquotedblleft{}The young man of the village was a hard worker.\textquotedblright{}
\end{quote}

\begin{cousinweb}
\textbf{\ja{はたらきもの}} (\emph{hatarakimono}) is a native Japanese word (wago), inherited rather than borrowed from Chinese, built from \emph{hataraki}, \textquotedblleft{}working\textquotedblright{} (from \emph{hataraku}), and \emph{mono}, \textquotedblleft{}person\textquotedblright{} (JMdict, KANJIDIC2).
\end{cousinweb}

\begin{grammarlens}[title={What you've built}]
One more word for telling what happened.
\end{grammarlens}

\subsection*{Your turn}
\begin{itemize}
\raggedright
  \item \emph{Say it:} \emph{hatarakimono}
  \item \emph{Say it:} \emph{hatarakimono}, once more
  \item \emph{Say it:} \emph{mura no wakai otoko wa hatarakimono deshita} — The young man of the village was a hard worker.
  \item \emph{Recall:} say the Japanese for to tell (a story), then the Japanese for a village, then say \emph{hatarakimono} again
\end{itemize}

\begin{tcolorbox}[breakable,colback=teal!4,colframe=teal!35!black,title={Before you move on}]
What does \ja{はたらきもの} mean? (\textquotedblleft{}A hard worker\textquotedblright{}.) Say it once more.
\end{tcolorbox}
